% GENERATED FILE. Edit canonical lessons, then run npm run generate:books.
% canonical-source-hash: fnv1a64:5b002f857395ec73
% canonical-lessons: SA-C45-morning, SA-C45-noon, SA-C45-evening, SA-C45-month, SA-C45-year

\chapter{Times of the Day}
\label{ch:sa-times-of-the-day}
\hlchaptermodality{45}

\begin{chapteropening}
\textbf{By the end of this chapter:} \emph{I can name morning, noon, evening, a month and a year in Sanskrit.}

Complete SA-C45-year and demonstrate the chapter promise.
\end{chapteropening}

\section[prātaḥ]{\sk{प्रातः} (prātaḥ) — morning}

\label{lesson:SA-C45-morning}



\begin{quote}
Before the new word: what did \emph{nīlaḥ} mean?
\end{quote}

\subsection*{You'll want to know: \sk{प्रातः}}
\textbf{\sk{प्रातः}} (\emph{prātaḥ}) — \textquotedblleft{}morning\textquotedblright{}.

\textbf{\sk{प्रातः}} is early morning: the stretch before the sun stands high, and the hour the first greeting of the day belongs to.

You have said \emph{pra-} at the front of a word twice already, in \textbf{\sk{प्रणामः}} and in \textbf{\sk{प्रसन्नः}}, where it carries the sense \textquotedblleft{}forward, out in front\textquotedblright{}. \textbf{\sk{प्रातः}} is the front end of the day.

Greek \emph{prōi}, \textquotedblleft{}early\textquotedblright{}, is the same small piece put to the same work, and Latin \emph{prae} is its cousin; English carries it as the \emph{pre-} in front of a thousand words.

\begin{grammarlens}[title={What you've built}]
The first part of the day, named.
\end{grammarlens}

\subsection*{Your turn}
Say these aloud:

\begin{itemize}
\raggedright
  \item \emph{prātaḥ}
  \item it once more, slowly
  \item \emph{nīlaḥ}, then \emph{prātaḥ}
\end{itemize}

\begin{tcolorbox}[breakable,colback=teal!4,colframe=teal!35!black,title={Before you move on}]
What does \sk{प्रातः} mean? (\textquotedblleft{}morning\textquotedblright{}.) Say it once more.
\end{tcolorbox}
\section[madhyāhnaḥ]{\sk{मध्याह्नः} (madhyāhnaḥ) — noon}

\label{lesson:SA-C45-noon}



\begin{quote}
Before the new word: what did \emph{prātaḥ} mean?
\end{quote}

\subsection*{You'll want to know: \sk{मध्याह्नः}}
\textbf{\sk{मध्याह्नः}} (\emph{madhyāhnaḥ}) — \textquotedblleft{}noon\textquotedblright{}.

\textbf{\sk{मध्याह्नः}} is midday, the sun at the top of its climb.

It is two pieces joined: \emph{madhya}, \textquotedblleft{}middle\textquotedblright{}, and \emph{ahan}, a day. Sanskrit keeps more than one word for a day, and this is the other one — \sk{दिनम्} is the day you count, \emph{ahan} the day you are standing in. English \emph{mid} is \emph{madhya}'s own worn-down cousin.

So noon is the middle of the day said in one breath, which is what noon is in most languages once you take the word apart.

\begin{grammarlens}[title={What you've built}]
The top of the day, between its two ends.
\end{grammarlens}

\subsection*{Your turn}
Say these aloud:

\begin{itemize}
\raggedright
  \item \emph{madhyāhnaḥ}
  \item it once more, slowly
  \item \emph{prātaḥ}, then \emph{madhyāhnaḥ}
\end{itemize}

\begin{tcolorbox}[breakable,colback=teal!4,colframe=teal!35!black,title={Before you move on}]
What does \sk{मध्याह्नः} mean? (\textquotedblleft{}noon\textquotedblright{}.) Say it once more.
\end{tcolorbox}
\section[sāyam]{\sk{सायम्} (sāyam) — evening}

\label{lesson:SA-C45-evening}



\begin{quote}
Before the new word: what did \emph{madhyāhnaḥ} mean?
\end{quote}

\subsection*{You'll want to know: \sk{सायम्}}
\textbf{\sk{सायम्}} (\emph{sāyam}) — \textquotedblleft{}evening\textquotedblright{}.

\textbf{\sk{सायम्}} is evening: the end of the working day, when a thing is set down rather than picked up.

Put it beside \sk{प्रातः} and you hold the two ends of the daylight — the greeting at one end and the leave-taking at the other. \emph{sāyaṁ prātaḥ}, \textquotedblleft{}evening and morning\textquotedblright{}, is one of the oldest paired phrases the language has.

It rests on a piece meaning \textquotedblleft{}to finish, to let go\textquotedblright{}, so evening is named as the slackening of the day rather than as its darkness.

\begin{grammarlens}[title={What you've built}]
Both ends of the daylight, and the middle between them.
\end{grammarlens}

\subsection*{Your turn}
Say these aloud:

\begin{itemize}
\raggedright
  \item \emph{sāyam}
  \item it once more, slowly
  \item \emph{madhyāhnaḥ}, then \emph{sāyam}
\end{itemize}

\begin{tcolorbox}[breakable,colback=teal!4,colframe=teal!35!black,title={Before you move on}]
What does \sk{सायम्} mean? (\textquotedblleft{}evening\textquotedblright{}.) Say it once more.
\end{tcolorbox}
\section[māsaḥ]{\sk{मासः} (māsaḥ) — a month}

\label{lesson:SA-C45-month}



\begin{quote}
Before the new word: what did \emph{sāyam} mean?
\end{quote}

\subsection*{You'll want to know: \sk{मासः}}
\textbf{\sk{मासः}} (\emph{māsaḥ}) — \textquotedblleft{}a month\textquotedblright{}.

\textbf{\sk{मासः}} is a month, and it is the moon's own word.

You have \sk{चन्द्रः} for the moon overhead. \textbf{\sk{मासः}} is that same moon counted: the stretch from one to the next.

Latin \emph{mēnsis}, Greek \emph{mēn}, English \emph{month} and \emph{moon} are all one word travelled west, and this is among the securest cousins in the whole family. Ends in \textbf{-\sk{अः}}.

\begin{grammarlens}[title={What you've built}]
A stretch of time measured by the moon you already know.
\end{grammarlens}

\subsection*{Your turn}
Say these aloud:

\begin{itemize}
\raggedright
  \item \emph{māsaḥ}
  \item it once more, slowly
  \item \emph{sāyam}, then \emph{māsaḥ}
\end{itemize}

\begin{tcolorbox}[breakable,colback=teal!4,colframe=teal!35!black,title={Before you move on}]
What does \sk{मासः} mean? (\textquotedblleft{}a month\textquotedblright{}.) Say it once more.
\end{tcolorbox}
\section[saṁvatsaraḥ]{\sk{संवत्सरः} (saṁvatsaraḥ) — a year}

\label{lesson:SA-C45-year}



\begin{quote}
Before the last word of the chapter: what did \emph{sāyam} mean, and what did \emph{māsaḥ} mean?
\end{quote}

\subsection*{You'll want to know: \sk{संवत्सरः}}
\textbf{\sk{संवत्सरः}} (\emph{saṁvatsaraḥ}) — \textquotedblleft{}a year\textquotedblright{}.

\textbf{\sk{संवत्सरः}} is a year: the whole round of months, come back to where it started.

The \emph{sam-} in front is the piece meaning \textquotedblleft{}together, complete\textquotedblright{} that also opens the name of the language itself. The \emph{vatsara} behind it is kin to Latin \emph{vetus}, \textquotedblleft{}old\textquotedblright{}, and to Greek \emph{etos}, a year — a year being what leaves a thing one year older. English \emph{veteran} sits on that same piece.

Beside \sk{मासः}, the small round, \textbf{\sk{संवत्सरः}} is the large one. Ends in \textbf{-\sk{अः}}.

\begin{grammarlens}[title={What you've built}]
Five stretches of time: morning, noon, evening, a month, a year.
\end{grammarlens}

\subsection*{Your turn}
Say these aloud:

\begin{itemize}
\raggedright
  \item \emph{saṁvatsaraḥ}
  \item it once more, slowly
  \item all five in order, then \emph{prātaḥ} and \emph{saṁvatsaraḥ} together
\end{itemize}

\begin{tcolorbox}[breakable,colback=teal!4,colframe=teal!35!black,title={Before you move on}]
What does \sk{संवत्सरः} mean? (\textquotedblleft{}a year\textquotedblright{}.) Say it once more.
\end{tcolorbox}
